\section{pandas}

\begin{frame}[fragile,t,label=cd-pandas]{2. Series e DataFrame}
\begin{itemize}
  \item \textbf{Series:} sequência de valores, cada um com seu índice.
  \item \textbf{DataFrame:} tabela com linhas e colunas que podem ter tipos diferentes.
\end{itemize}
\begin{columns}[T,onlytextwidth]
\begin{column}{0.54\textwidth}
\EMASCcodelabel{Exemplo prático}
\begin{lstlisting}
serie = pd.Series([5.9, 6.7, 7.5])
exemplo = pd.DataFrame({
    "id": [1, 2, 3],
    "nota": [5.9, 6.7, 7.5]
})
print(serie)
print(exemplo)
\end{lstlisting}
\end{column}
\begin{column}{0.42\textwidth}
\EMASCcodelabel{DataFrame}
\begin{tabular}{rrr}
\toprule
\textbf{Índice} & \textbf{id} & \textbf{nota} \\
\midrule
0 & 1 & 5.9 \\
1 & 2 & 6.7 \\
2 & 3 & 7.5 \\
\bottomrule
\end{tabular}
\smallskip

A coluna \texttt{nota} com seus índices é uma Series.
\end{column}
\end{columns}
\end{frame}

\begin{frame}[fragile,t]{2.1. Leitura de uma base de dados}
CSV é um arquivo de texto com dados separados por um delimitador, geralmente vírgula.
\EMASCcodelabel{Sintaxe}
\texttt{pd.read\_csv("notas.csv")}
\EMASCcodelabel{Leitura do arquivo de notas}
\begin{lstlisting}
df = pd.read_csv(fonte)
print(df.head())
\end{lstlisting}
\begin{EMASCnote}
A função aceita um caminho local ou uma URL. A variável \texttt{fonte} é definida na célula de configuração do notebook.
\end{EMASCnote}
\medskip
Para arquivos separados por ponto e vírgula, utilize \texttt{sep=";"}.
\end{frame}

\begin{frame}[fragile,t]{2.2. Inspeção dos dados}
Antes de realizar cálculos, é importante conhecer o formato da tabela e o significado das colunas.
\EMASCcodelabel{Exemplo prático}
\begin{lstlisting}
print("Primeiras linhas:\n", df.head())
print("Ultimas linhas:\n", df.tail())
print("Linhas e colunas:", df.shape)
print("Colunas:", df.columns)
df.info()   # Tipos e quantidade de valores presentes
print(df[["horas_estudo", "frequencia", "nota"]].describe())
\end{lstlisting}
\begin{EMASCnote}
\texttt{shape} retorna \texttt{(18, 5)}. \texttt{describe()} resume as variáveis numéricas; o ID foi excluído por ser um identificador.
\end{EMASCnote}
\end{frame}

\begin{frame}[fragile,t]{2.3. Seleção de linhas e colunas}
\texttt{loc} seleciona por rótulos. \texttt{iloc} seleciona por posições inteiras, iniciadas em zero.
\EMASCcodelabel{Exemplo prático}
\begin{lstlisting}
print(df["nota"])                      # Uma Series
print(df[["curso", "nota"]])           # Um DataFrame
print(df.loc[0:2, ["id", "nota"]])     # Rotulos
print(df.iloc[0:3, [0, 4]])            # Posicoes
\end{lstlisting}
\begin{EMASCnote}
Os dois últimos comandos retornam as mesmas linhas e colunas. No \texttt{loc}, inclui-se o rótulo final; no \texttt{iloc}, exclui-se a posição final.
\end{EMASCnote}
\end{frame}

\begin{frame}[fragile,t]{2.4. Filtros e ordenação}
Filtros selecionam linhas que satisfazem condições. \texttt{isna()} identifica valores ausentes.
\EMASCcodelabel{Exemplo prático}
\begin{lstlisting}
print(df[df["curso"] == "Matematica"])
alto = (df["nota"] >= 7) & (df["frequencia"] >= 0.8)
print(df[alto])
atencao = (df["nota"] < 7) | (df["nota"].isna())
print(df[atencao])
print(df.sort_values("nota", ascending=False).head())
\end{lstlisting}
\begin{EMASCnote}
Use \texttt{\&}, \texttt{|} e \texttt{\textasciitilde} no lugar de \texttt{and}, \texttt{or} e \texttt{not} para comparar linha a linha. \texttt{ascending=False} ordena do maior para o menor.
\end{EMASCnote}
\end{frame}
